% GENERATED FILE. Edit canonical lessons, then run npm run generate:books.
% canonical-source-hash: fnv1a64:d9d61a14226a6ad6
% canonical-lessons: PT-C285-entrevistar, PT-C285-entrevistado, PT-C285-ouvinte, PT-C285-telespectador, PT-C285-leitor

\chapter{Quem ouve e quem lê}
\label{ch:pt-quem-ouve-e-quem-le}
\hlchaptermodality{285}

\begin{chapteropening}
\textbf{By the end of this chapter:} \emph{I can talk about interviews and about the people who listen to, watch and read the news.}

Complete the last lesson of chapter 285: I can talk about interviews and about the people who listen to, watch and read the news.
\end{chapteropening}

\section[entrevistar]{entrevistar — to interview}

\label{lesson:PT-C285-entrevistar}



\begin{quote}
Before the new one: say the Portuguese for the newspaper piece, then the Portuguese for the columnist.
\end{quote}

\subsection*{You'll want to know: entrevistar}
\textbf{entrevistar} — \textquotedblleft{}to interview\textquotedblright{}.

\begin{quote}
A jornalista entrevistou o novo treinador.
\end{quote}

\begin{quote}
\emph{The journalist interviewed the new coach.}
\end{quote}

Built on \emph{entrevista}.

\begin{cousinweb}
\emph{entrevistar} is built on \emph{entrevista} \textquotedblleft{}interview\textquotedblright{}, which is usually said to follow French \emph{entrevue} (EtymDB has no source for it); \emph{vista} goes back to Latin \emph{vidēre} \textquotedblleft{}to see\textquotedblright{} (Lewis \& Short).
\end{cousinweb}

\begin{grammarlens}[title={What you've built}]
One more word for telling, explaining and describing at B1.
\end{grammarlens}

\subsection*{Your turn}
\begin{itemize}
\raggedright
  \item \emph{Say it:} \emph{entrevistar}
  \item \emph{Say it:} \emph{entrevistar}, once more
  \item \emph{Say it:} a sentence of your own with \emph{entrevistar}
  \item \emph{Recall:} say the Portuguese for the newspaper piece, then the Portuguese for the columnist, then say \emph{entrevistar} again
\end{itemize}

\begin{tcolorbox}[breakable,colback=teal!4,colframe=teal!35!black,title={Before you move on}]
What does entrevistar mean? (\textquotedblleft{}To interview\textquotedblright{}.) Say it once more.
\end{tcolorbox}
\section[o entrevistado]{o entrevistado — the interviewee}

\label{lesson:PT-C285-entrevistado}



\begin{quote}
Before the new one: say the Portuguese for the columnist, then the Portuguese for to interview.
\end{quote}

\subsection*{You'll want to know: o entrevistado}
\textbf{o entrevistado} — \textquotedblleft{}the interviewee\textquotedblright{}.

\begin{quote}
O entrevistado respondeu com calma.
\end{quote}

\begin{quote}
\emph{The interviewee answered calmly.}
\end{quote}

A woman is \emph{a entrevistada}.

\begin{cousinweb}
\emph{entrevistado} is the past participle of \emph{entrevistar} used as a noun; \emph{vista} goes back to Latin \emph{vidēre} \textquotedblleft{}to see\textquotedblright{} (Lewis \& Short).
\end{cousinweb}

\begin{grammarlens}[title={What you've built}]
One more word for telling, explaining and describing at B1.
\end{grammarlens}

\subsection*{Your turn}
\begin{itemize}
\raggedright
  \item \emph{Say it:} \emph{o entrevistado}
  \item \emph{Say it:} \emph{o entrevistado}, once more
  \item \emph{Say it:} a sentence of your own with \emph{o entrevistado}
  \item \emph{Recall:} say the Portuguese for the columnist, then the Portuguese for to interview, then say \emph{o entrevistado} again
\end{itemize}

\begin{tcolorbox}[breakable,colback=teal!4,colframe=teal!35!black,title={Before you move on}]
What does o entrevistado mean? (\textquotedblleft{}The interviewee\textquotedblright{}.) Say it once more.
\end{tcolorbox}
\section[o ouvinte]{o ouvinte — the listener}

\label{lesson:PT-C285-ouvinte}



\begin{quote}
Before the new one: say the Portuguese for to interview, then the Portuguese for the interviewee.
\end{quote}

\subsection*{You'll want to know: o ouvinte}
\textbf{o ouvinte} — \textquotedblleft{}the listener\textquotedblright{}.

\begin{quote}
Um ouvinte telefonou para a rádio.
\end{quote}

\begin{quote}
\emph{A listener phoned the radio station.}
\end{quote}

The same form for a woman: \emph{a ouvinte}.

\begin{cousinweb}
\emph{ouvinte} is inherited from Latin \emph{audiēns} \textquotedblleft{}hearing, a hearer\textquotedblright{} (EtymDB), from \emph{audīre} \textquotedblleft{}to hear\textquotedblright{}.
\end{cousinweb}

\begin{grammarlens}[title={What you've built}]
One more word for telling, explaining and describing at B1.
\end{grammarlens}

\subsection*{Your turn}
\begin{itemize}
\raggedright
  \item \emph{Say it:} \emph{o ouvinte}
  \item \emph{Say it:} \emph{o ouvinte}, once more
  \item \emph{Say it:} a sentence of your own with \emph{o ouvinte}
  \item \emph{Recall:} say the Portuguese for to interview, then the Portuguese for the interviewee, then say \emph{o ouvinte} again
\end{itemize}

\begin{tcolorbox}[breakable,colback=teal!4,colframe=teal!35!black,title={Before you move on}]
What does o ouvinte mean? (\textquotedblleft{}The listener\textquotedblright{}.) Say it once more.
\end{tcolorbox}
\section[o telespectador]{o telespectador — the TV viewer}

\label{lesson:PT-C285-telespectador}



\begin{quote}
Before the new one: say the Portuguese for the interviewee, then the Portuguese for the listener.
\end{quote}

\subsection*{You'll want to know: o telespectador}
\textbf{o telespectador} — \textquotedblleft{}the TV viewer\textquotedblright{}.

\begin{quote}
Cada telespectador pode votar por telefone.
\end{quote}

\begin{quote}
\emph{Each viewer can vote by phone.}
\end{quote}

Built on \emph{tele-} and \emph{espectador}, ‘onlooker’.

\begin{cousinweb}
\emph{telespectador} joins \emph{tele-}, from Greek \emph{tēle} \textquotedblleft{}far off\textquotedblright{}, and \emph{espectador}, from Latin \emph{spectātor} \textquotedblleft{}an onlooker\textquotedblright{}, from \emph{spectāre} \textquotedblleft{}to watch\textquotedblright{} (Lewis \& Short).
\end{cousinweb}

\begin{grammarlens}[title={What you've built}]
One more word for telling, explaining and describing at B1.
\end{grammarlens}

\subsection*{Your turn}
\begin{itemize}
\raggedright
  \item \emph{Say it:} \emph{o telespectador}
  \item \emph{Say it:} \emph{o telespectador}, once more
  \item \emph{Say it:} a sentence of your own with \emph{o telespectador}
  \item \emph{Recall:} say the Portuguese for the interviewee, then the Portuguese for the listener, then say \emph{o telespectador} again
\end{itemize}

\begin{tcolorbox}[breakable,colback=teal!4,colframe=teal!35!black,title={Before you move on}]
What does o telespectador mean? (\textquotedblleft{}The TV viewer\textquotedblright{}.) Say it once more.
\end{tcolorbox}
\section[o leitor]{o leitor — the reader}

\label{lesson:PT-C285-leitor}



\begin{quote}
Before the new one: say the Portuguese for the listener, then the Portuguese for the TV viewer.
\end{quote}

\subsection*{You'll want to know: o leitor}
\textbf{o leitor} — \textquotedblleft{}the reader\textquotedblright{}.

\begin{quote}
Um leitor escreveu uma carta ao jornal.
\end{quote}

\begin{quote}
\emph{A reader wrote a letter to the paper.}
\end{quote}

A woman is \emph{a leitora}. Also a device: \emph{um leitor de CD}.

\begin{cousinweb}
\emph{leitor} comes from Latin \emph{lēctor} \textquotedblleft{}a reader\textquotedblright{} (EtymDB), from \emph{legere} \textquotedblleft{}to read\textquotedblright{}.
\end{cousinweb}

\begin{grammarlens}[title={What you've built}]
One more word for telling, explaining and describing at B1.
\end{grammarlens}

\subsection*{Your turn}
\begin{itemize}
\raggedright
  \item \emph{Say it:} \emph{o leitor}
  \item \emph{Say it:} \emph{o leitor}, once more
  \item \emph{Say it:} a sentence of your own with \emph{o leitor}
  \item \emph{Recall:} say the Portuguese for the listener, then the Portuguese for the TV viewer, then say \emph{o leitor} again
\end{itemize}

\begin{tcolorbox}[breakable,colback=teal!4,colframe=teal!35!black,title={Before you move on}]
What does o leitor mean? (\textquotedblleft{}The reader\textquotedblright{}.) Say it once more.
\end{tcolorbox}
